%% ===========================================================================
%%  frontmatter/kata-pengantar.tex — Kata Pengantar
%% ===========================================================================
\babfrontmatter{Kata Pengantar}

\begin{singlespace}
Puji syukur penulis sampaikan kepada Tuhan Yang Maha Esa karena laporan magang
ini dapat diselesaikan. Laporan disusun sebagai pertanggungjawaban pelaksanaan
magang pada Program Studi \ProgramStudi{} di \NamaKampus. Kegiatan magang
dilaksanakan di \InstansiMagang{} pada periode \PeriodeMagang. Penyusunan
laporan ini tidak terlepas dari bantuan berbagai pihak yang telah memberikan
arahan, kesempatan, dan dukungan kepada penulis.

Penulis mengucapkan terima kasih kepada \NamaPembimbing{} sebagai dosen
pembimbing yang telah memberikan arahan mengenai isi dan format laporan.
Ucapan terima kasih juga penulis sampaikan kepada \NamaKetuaProgramStudi{}
sebagai Ketua Program Studi \ProgramStudi{} atas dukungan administratif
selama kegiatan magang. Penulis berterima kasih kepada pimpinan \NamaKampus{}
atas kesempatan mengikuti program magang ini. Dukungan tersebut membantu
penulis memenuhi kewajiban akademik dan menyelesaikan laporan tepat pada
waktunya.

Penulis menyampaikan terima kasih kepada Ian Joseph Chandra sebagai mentor di
Broadband Multimedia Wireless Laboratory atas bimbingan selama pelaksanaan
pekerjaan. Penulis juga berterima kasih kepada dosen, peneliti, dan anggota
laboratorium yang telah memberikan masukan mengenai kebutuhan sistem. Diskusi
bersama mereka membantu penulis memahami konteks pekerjaan penelitian. Ucapan
terima kasih juga disampaikan kepada rekan mahasiswa magang yang telah
berbagi pengalaman dan bantuan teknis.

Penulis mengucapkan terima kasih kepada orang tua dan keluarga atas doa serta
dukungan yang diberikan selama masa studi dan magang. Dukungan tersebut
menjadi bagian penting dalam penyelesaian laporan ini. Penulis menyadari bahwa
laporan masih memiliki keterbatasan, terutama karena sistem yang dibahas baru
diuji pada sejumlah skenario. Karena itu, penulis mengharapkan kritik dan
saran yang membangun untuk perbaikan pada masa berikutnya. Semoga laporan ini
memberi manfaat bagi pembaca dan bagi pengembangan pekerjaan serupa di
laboratorium.

\vspace{2\baselineskip}
\begin{flushright}
\Kota, \TanggalKataPengantar\\[2\baselineskip]
Penulis
\end{flushright}
\end{singlespace}

